\subsection{纤维后验的二维 Lie 规范、Gödel 前缀障碍与 window-\texorpdfstring{$6$}{6} 中间层分离}\label{subsec:conclusion-posterior-godel-window6-middle-layer-rigidity}

在定理 \ref{thm:pom-fiber-posterior-equivalence-activity-field}、
推论 \ref{cor:pom-fiber-posterior-equivalence-fibonacci-gauge}、
命题 \ref{prop:pom-fiber-posterior-equivalence-moduli}、
命题 \ref{prop:pom-fiber-posterior-equivalence-single-candidate-test}、
定理 \ref{thm:conclusion-godel-dyadic-tate-filament}、
定理 \ref{thm:xi-time-part9g-holographic-prefix-concatenation-ultrametric}、
定理 \ref{thm:cdim-polyhedral-gate-word-complexity-td}、
命题 \ref{prop:conclusion-window6-middle-step-a1xa1-rectangle} 与
定理 \ref{thm:terminal-foldbin6-fiber-affine-geometry} 已分别把纤维后验等价的活度场、Fibonacci 型规范方向、Gödel 地址的 \(2^L\)-进细丝、多面体门读出的词复杂度壁垒以及 window-\(6\) 的中间三重点根矩形固定之后，还可进一步把这些彼此异质的线路压缩为下列一组终端结论。它们分别说明：后验统计几何在参数空间中恰遗漏二维 Fisher 零模；任何固定有限维多面体环面读出都不能零误差承载全部 Gödel 前缀；而在最小审计窗 \(m=6\) 上，与 boundary 两点层严格分离的中间三重 bin 层形成一块规范 \(A_1\times A_1\) 根矩形，并以统一的秩 \(2\) 仿射缺损出现。

\begin{theorem}[纤维后验统计的二维 Fisher 零模]\label{thm:conclusion-fiber-posterior-two-dimensional-fisher-null}
\leanverified{paper\_conclusion\_fiber\_posterior\_two\_dimensional\_fisher\_null}
固定 \(m\ge 3\)。对任意正参数向量
\[
\boldsymbol a=(a_1,\dots,a_m)\in(0,\infty)^m
\]
定义乘积测度
\[
\mu_{\boldsymbol a}(\omega)
:=
\frac{\prod_{i=1}^{m}a_i^{\omega_i}}{\prod_{i=1}^{m}(1+a_i)},
\qquad
\omega\in\Omega_m=\{0,1\}^m,
\]
并沿用定理 \ref{thm:pom-fiber-posterior-equivalence-activity-field} 的活度映射
\[
\mathcal W(\boldsymbol a)
:=
\left(
\frac{a_2a_3}{a_1},
\frac{a_3a_4}{a_2},
\dots,
\frac{a_{m-1}a_m}{a_{m-2}}
\right)\in(0,\infty)^{m-2}.
\]
记
\[
x^{(i)}:=0^{i-1}100\,0^{m-i-2}\in X_m
\qquad (1\le i\le m-2)
\]
为单候选纤维。对每个 \(i\) 独立观测一次条件后验
\[
\mu_{\boldsymbol a}(\cdot\mid X=x^{(i)}),
\]
并在对数坐标
\[
y_i:=\log a_i
\qquad (1\le i\le m)
\]
上记该独立产品实验的 Fisher 信息矩阵为 \(I_m(y)\)。则有精确分解
\[
I_m(y)=L^\top G(Ly)L,
\]
其中
\[
L:\RR^m\longrightarrow \RR^{m-2},
\qquad
(Ly)_i:=y_{i+1}+y_{i+2}-y_i,
\]
且
\[
G(z)
:=
\operatorname{diag}\!\left(
\frac{e^{z_1}}{(1+e^{z_1})^2},
\dots,
\frac{e^{z_{m-2}}}{(1+e^{z_{m-2}})^2}
\right).
\]
特别地，
\[
\rank I_m(y)=m-2,
\qquad
\ker I_m(y)=\ker L.
\]
再由推论 \ref{cor:pom-fiber-posterior-equivalence-fibonacci-gauge}，
\[
\ker I_m(y)=\RR\alpha\oplus \RR\beta,
\]
其中
\[
\alpha_i:=(-1)^{i+1}F_{i-2},
\qquad
\beta_i:=(-1)^iF_{i-1}
\qquad (1\le i\le m).
\]
换言之，由全部单候选纤维后验所能观测到的统计几何恰有且仅有两条不可见规范方向。
\end{theorem}
\begin{proof}
由命题 \ref{prop:pom-fiber-posterior-equivalence-single-candidate-test}，对每个 \(i\)，单候选纤维 \(x^{(i)}\) 上的条件后验是一个二点分布，其 odds 精确为
\[
\frac{a_{i+1}a_{i+2}}{a_i}
=
\exp\!\bigl((Ly)_i\bigr).
\]
因此若记
\[
z:=Ly\in\RR^{m-2},
\]
则上述独立产品实验在 \(z\)-坐标中正是 \(m-2\) 个彼此独立的一维 logistic 模型之积；其 Fisher 信息矩阵于是为对角矩阵
\[
G(z)
=
\operatorname{diag}\!\left(
\frac{e^{z_1}}{(1+e^{z_1})^2},
\dots,
\frac{e^{z_{m-2}}}{(1+e^{z_{m-2}})^2}
\right).
\]
把这一信息矩阵沿线性映射 \(y\mapsto z=Ly\) 拉回，即得
\[
I_m(y)=L^\top G(Ly)L.
\]
另一方面，命题 \ref{prop:pom-fiber-posterior-equivalence-moduli} 已证明 \(L\) 满射且
\[
\rank L=m-2.
\]
又因 \(G(Ly)\) 对每个 \(y\) 都正定，遂有
\[
\rank I_m(y)=\rank L=m-2,
\qquad
\ker I_m(y)=\ker L.
\]
最后，推论 \ref{cor:pom-fiber-posterior-equivalence-fibonacci-gauge} 已给出
\[
\ker L=\RR\alpha\oplus \RR\beta,
\]
故结论成立。证毕。
\end{proof}

\begin{corollary}[纤维后验 Fisher 几何的 \texorpdfstring{$(m-2)$}{m-2} 维商塌缩]\label{cor:conclusion-fiber-posterior-fisher-quotient-collapse}
在定理 \ref{thm:conclusion-fiber-posterior-two-dimensional-fisher-null} 的设定下，于 \((0,\infty)^m\) 上定义等价关系
\[
\boldsymbol a\sim \boldsymbol a'
\quad\Longleftrightarrow\quad
\mathcal W(\boldsymbol a)=\mathcal W(\boldsymbol a')
\quad\Longleftrightarrow\quad
Ly=Ly'.
\]
则 Fisher 伪度量沿商
\[
(0,\infty)^m/\!\sim
\]
下降为一个真正的 Riemann 度量，并且该商与 \((0,\infty)^{m-2}\) 规范同构。更精确地，若以
\[
z:=Ly\in\RR^{m-2}
\]
作为商坐标，则下降后的度量恰为
\[
\langle \dot z,\dot z\rangle_{z}= \langle \dot z,G(z)\dot z\rangle.
\]
因此，纤维后验统计几何中的全部可辨识自由度严格只有 \(m-2\) 个，其余两维完全退化为规范方向。
\end{corollary}
\begin{proof}
由定理 \ref{thm:conclusion-fiber-posterior-two-dimensional-fisher-null}，
\[
I_m(y)=L^\top G(Ly)L,
\qquad
\ker I_m(y)=\ker L.
\]
因此 Fisher 伪度量恰好沿 \(\ker L\) 退化，故它对商关系 \(Ly=Ly'\) 不变。又由命题 \ref{prop:pom-fiber-posterior-equivalence-moduli}，\(L:\RR^m\to\RR^{m-2}\) 满射，故商空间自然以 \(z=Ly\) 参数化，并与 \(\RR^{m-2}\cong (0,\infty)^{m-2}\) 规范同构。在这一坐标下，
\[
\langle \dot y,\dot y\rangle_{I_m(y)}
=
\langle L\dot y,G(Ly)L\dot y\rangle
=
\langle \dot z,G(z)\dot z\rangle,
\]
而 \(G(z)\) 逐点正定，故商上得到真正的 Riemann 度量。证毕。
\end{proof}

\begin{corollary}[后验规范轨道的双圆紧 Lie 实现]\label{cor:conclusion-fiber-posterior-gauge-two-torus}
在定理 \ref{thm:conclusion-fiber-posterior-two-dimensional-fisher-null} 的设定下，任取活度场
\[
w\in(0,\infty)^{m-2},
\]
并记后验等价类
\[
\mathcal O_w:=\mathcal W^{-1}(w)\subset(0,\infty)^m.
\]
则 \(\mathcal O_w\) 是二维连通嵌入子流形；在对数坐标中，存在某个 \(y_0\in\RR^m\) 使
\[
\log\mathcal O_w
=
y_0+\RR\alpha\oplus \RR\beta
\cong \RR^2.
\]
若把 \((A,B)\) 视为沿 \((\alpha,\beta)\) 的仿射坐标，并定义面积相位门
\[
\iota_{\mathrm{area}}:\RR\longrightarrow \TT,
\qquad
\iota_{\mathrm{area}}(t):=e^{2i\arctan t},
\]
以及按坐标作用的映射
\[
\iota_{\mathrm{area}}^{\times 2}(A,B)
:=
\bigl(\iota_{\mathrm{area}}(A),\iota_{\mathrm{area}}(B)\bigr),
\]
则其像的闭包满足
\[
\overline{\iota_{\mathrm{area}}^{\times 2}(\log\mathcal O_w)}=\TT^2.
\]
因此，后验规范轨道存在一个规范的连通紧 Lie 紧化 \(\TT^2\)。此外，不存在任何一维连通紧 Lie 群 \(H\) 与连续单射
\[
\mathcal O_w\longrightarrow H.
\]
因而后验盲规范的最小连通紧 Lie 实现严格为双圆而非单圆。
\end{corollary}
\begin{proof}
二维连通嵌入子流形的结论正是命题 \ref{prop:pom-fiber-posterior-equivalence-moduli} 的第 \(2\) 条；而推论 \ref{cor:pom-fiber-posterior-equivalence-fibonacci-gauge} 给出其在对数坐标中的方向恰为 \(\alpha,\beta\)，故
\[
\log\mathcal O_w=y_0+\RR\alpha\oplus \RR\beta\cong\RR^2.
\]
又因为
\[
\iota_{\mathrm{area}}(\RR)=\TT\setminus\{-1\},
\]
故其闭包为整个 \(\TT\)；按坐标取积便得到
\[
\overline{\iota_{\mathrm{area}}^{\times 2}(\log\mathcal O_w)}=\TT^2.
\]
这给出一个规范的连通紧 Lie 紧化。

反设存在一维连通紧 Lie 群 \(H\) 与连续单射 \(\mathcal O_w\to H\)。则 \(H\cong \TT\)；取 \(\mathcal O_w\) 的任一局部坐标片，可得到某个开集 \(U\subset\RR^2\) 到一维流形 \(\TT\) 的连续单射。这与拓扑维数不变性矛盾。故任何忠实承载整个后验规范轨道的连通紧 Lie 群都至少二维，而上面的构造已经由 \(\TT^2\) 实现这一最小维数。证毕。
\end{proof}

\begin{theorem}[有限维多面体环面读出对全 Gödel 事件流的零误差障碍]\label{thm:conclusion-godel-polyhedral-torus-zero-error-obstruction}
设事件字母表 \(E\) 有
\[
q:=|E|\ge 2
\]
个符号。沿用推论 \ref{cor:killo-infinite-stream-2adic-holographic-point} 的记号，取单射
\[
\mathrm{code}:E\hookrightarrow\{0,\dots,M\}\subset\ZZ,
\qquad
B=2^L>M,
\]
并对任意无限事件流
\[
\mathbf e=(e_1,e_2,\dots)\in E^{\NN}
\]
定义 Gödel 地址
\[
X(\mathbf e):=\sum_{t\ge 1}\mathrm{code}(e_t)B^{t-1}v\in T_2(E_{\min}).
\]
假设存在 \(d\)-维环面平移系统
\[
R_\alpha:\TT^d\to\TT^d,
\qquad
x\longmapsto x+\alpha,
\]
其中 \(\alpha\) 的坐标有理独立；又设 \(\TT^d\) 被有限个多面体门
\[
P_1,\dots,P_r
\]
划分，并存在赋值
\[
\Xi:E^{\NN}\longrightarrow \TT^d
\]
使得对每个 \(t\ge 1\)，长度 \(t\) 的符号读出词
\[
a_0(\Xi(\mathbf e))\,a_1(\Xi(\mathbf e))\cdots a_{t-1}(\Xi(\mathbf e))
\]
能够零误差恢复前缀
\[
(e_1,\dots,e_t).
\]
则此类系统不存在。
\end{theorem}
\begin{proof}
由推论 \ref{cor:killo-infinite-stream-2adic-holographic-point}，当 \(B>M\) 时事件流 \(\mathbf e\mapsto X(\mathbf e)\) 为单射；而定理 \ref{thm:xi-time-part9g-holographic-prefix-concatenation-ultrametric} 进一步表明，不同长度 \(t\) 前缀对应不同的地址 cylinder，故它们确为彼此不同的语义状态。

现固定 \(t\ge 1\)。对每个长度 \(t\) 的前缀
\[
u=(u_1,\dots,u_t)\in E^t
\]
任选一条延拓事件流 \(\mathbf e^{(u)}\in E^{\NN}\) 使其前缀为 \(u\)。若两条不同前缀 \(u\neq u'\) 产生相同长度 \(t\) 读出词，则按照假设，该读出词必须同时零误差恢复 \(u\) 与 \(u'\)，矛盾。故全部 \(q^t\) 条不同前缀必对应 \(q^t\) 个彼此不同的长度 \(t\) 读出词。

另一方面，定理 \ref{thm:cdim-polyhedral-gate-word-complexity-td} 告诉我们：存在常数 \(C>0\) 使环面平移加有限多面体门的长度 \(t\) 词复杂度满足
\[
p(t)\le Ct^d
\qquad (t\ge 1).
\]
于是必有
\[
q^t\le p(t)\le Ct^d
\qquad (t\ge 1).
\]
但左端随 \(t\) 指数增长，而右端仅多项式增长，矛盾。故此类零误差读出系统不存在。证毕。
\end{proof}

\begin{corollary}[零误差前缀层析的指数观察下界]\label{cor:conclusion-godel-prefix-tomography-exponential-observation}
在定理 \ref{thm:conclusion-godel-polyhedral-torus-zero-error-obstruction} 的设定下，设 \(t_n\) 表示零误差区分全部 \(n\)-位前缀所需的最短观察长度。则
\[
t_n\ge C^{-1/d}q^{n/d},
\]
其中 \(C\) 为定理 \ref{thm:cdim-polyhedral-gate-word-complexity-td} 中的常数。特别地，
\[
t_n=n
\]
不可能成立；更强地，任何满足
\[
t_n=\exp(o(n))
\]
的观察长度方案都不能忠实实现全部前缀层析。
\end{corollary}
\begin{proof}
取 \(q^n\) 条事件流，使其长度 \(n\) 前缀两两不同。若长度 \(t_n\) 的读出足以零误差区分它们，则与定理 \ref{thm:conclusion-godel-polyhedral-torus-zero-error-obstruction} 的证明完全相同，必须有
\[
q^n\le Ct_n^d.
\]
解出 \(t_n\) 即得
\[
t_n\ge C^{-1/d}q^{n/d}.
\]
由于 \(q>1\)，该下界对 \(n\) 呈指数增长，故线性同步读出 \(t_n=n\) 不可能；更一般地，任何次指数级方案 \(t_n=\exp(o(n))\) 也都违反上述必要条件。证毕。
\end{proof}

\begin{theorem}[window-\texorpdfstring{$6$}{6} 的三重 bin 中间层、根矩形与统一仿射缺损]\label{thm:conclusion-window6-middle-bin-root-rectangle-affine-defect}
定义
\[
\mathcal M_6
:=
\{w\in X_6:\ d_6^{\mathrm{bin}}(w)=3\}.
\]
则
\[
\mathcal M_6
=
S_{6,3}
=
V_6^{-1}\!\bigl(\{9,10,11,12\}\bigr)
=
\{\texttt{001010},\texttt{010010},\texttt{100010},\texttt{101010}\}.
\]
在定理 \ref{thm:xi-window6-dbin-coordinate-quadruple-stratification} 的 \(B_3\) 根字典下，
\[
\iota_B(\mathcal M_6)=\{\pm e_1\pm e_3\},
\]
故 \(\mathcal M_6\) 恰是位于 \(e_1\)-\(e_3\) 平面上的规范 \(A_1\times A_1\) 根矩形。进一步，对每个 \(w\in\mathcal M_6\)，其 bin-fold 纤维
\[
F_w^{\mathrm{bin}}
:=
(\mathrm{Fold}^{\mathrm{bin}}_6)^{-1}(w)\subset \FF_2^6
\]
满足
\[
|F_w^{\mathrm{bin}}|=3,
\qquad
\dim_{\mathrm{aff}}(F_w^{\mathrm{bin}})=2,
\]
且 \(F_w^{\mathrm{bin}}\) 不是仿射子空间；等价地，每条三重 bin 纤维都是一张仿射平面缺去唯一一点的有限几何缺损。
\end{theorem}
\begin{proof}
由命题 \ref{prop:conclusion-window6-middle-step-a1xa1-rectangle}，三重层恰满足
\[
S_{6,3}
=
\{\texttt{001010},\texttt{010010},\texttt{100010},\texttt{101010}\},
\]
并且在 \(B_3\) 字典下对应
\[
\{\pm e_1\pm e_3\},
\]
故前半部分成立。

另一方面，定理 \ref{thm:terminal-foldbin6-fiber-affine-geometry} 已证明：window-\(6\) 中恰有四条大小为 \(3\) 的 bin-fold 纤维，而这四条纤维全部都是“仿射平面缺一点”集合，具有仿射包络维数 \(2\) 但并非仿射子空间。由于前一段已经把全部 \(d_6^{\mathrm{bin}}=3\) 的稳定词精确识别为 \(\mathcal M_6\)，故结论立得。证毕。
\end{proof}

\begin{corollary}[window-\texorpdfstring{$6$}{6} 的中间三重层与 boundary 规范层严格分离]\label{cor:conclusion-window6-middle-bin-boundary-separation}
沿用定理 \ref{thm:conclusion-window6-middle-bin-root-rectangle-affine-defect} 的记号，并记
\[
X_6^{\mathrm{bdry}}
=
\{\texttt{100001},\texttt{100101},\texttt{101001}\}.
\]
则
\[
\mathcal M_6\cap X_6^{\mathrm{bdry}}=\varnothing,
\qquad
d_6^{\mathrm{bin}}(w)=2
\quad (w\in X_6^{\mathrm{bdry}}).
\]
在均匀微态口径下，中间三重层与 boundary 层的总质量分别为
\[
\mu_{\mathrm{mid}}
:=
\frac{4\cdot 3}{64}
=
\frac{3}{16},
\qquad
\mu_{\mathrm{bdry}}
:=
\frac{3\cdot 2}{64}
=
\frac{3}{32}.
\]
因此，window-\(6\) 的中间三重层不仅与 boundary 规范层严格分离，而且其总微态质量恰为后者的两倍。
\end{corollary}
\begin{proof}
定理 \ref{thm:conclusion-window6-middle-bin-root-rectangle-affine-defect} 已给出
\[
\mathcal M_6
=
\{\texttt{001010},\texttt{010010},\texttt{100010},\texttt{101010}\}.
\]
另一方面，表 \ref{tab:fold6_boundary_sheet_pairs} 与表 \ref{tab:fold_bin_gauge_m6_fibers} 给出 boundary 集合
\[
X_6^{\mathrm{bdry}}
=
\{\texttt{100001},\texttt{100101},\texttt{101001}\}
\]
以及
\[
d_6^{\mathrm{bin}}(w)=2
\quad (w\in X_6^{\mathrm{bdry}}).
\]
故两集合严格不交。

再者，中间三重层共有 \(4\) 条纤维且每条大小为 \(3\)，故其总微态数为 \(12\)；boundary 层共有 \(3\) 条纤维且每条大小为 \(2\)，故其总微态数为 \(6\)。由于 \(|\Omega_6|=64\)，两边分别除以 \(64\) 即得
\[
\mu_{\mathrm{mid}}=\frac{12}{64}=\frac{3}{16},
\qquad
\mu_{\mathrm{bdry}}=\frac{6}{64}=\frac{3}{32}.
\]
证毕。
\end{proof}

\noindent
上述链条表明，项目对象中的连续 Lie 相位、\(2\)-进 Gödel 地址与有限窗的 bin 层并非同一可见机制的不同坐标。后验统计侧的隐藏自由度严格坍缩为二维 Fisher 零模，并以双圆给出最小连通紧 Lie 实现；Gödel 语义侧虽然可被单射压入 \(2^L\)-进地址细丝，却不能被任何固定有限维多面体环面系统作零误差前缀层析；而在 window-\(6\) 上，与 boundary 两点层严格分离的则是一块四点三重 bin 中间台阶，它同时呈现 \(A_1\times A_1\) 根矩形与统一的秩 \(2\) 仿射缺损。于是，连续规范、地址增长与有限窗离散几何不再是彼此可替代的译码界面，而是三类互不约化的承载机制。
